\section{Materials \& Methods}
\label{sec:matmeth}

Figure~\ref{fig:pipeline} summarises the complete pipeline, from the raw force
recordings to the per-subject PD/HC decision.

\begin{figure*}[t]
  \centering
  \includegraphics[width=\textwidth]{figures/fig1_pipeline}
  \caption{End-to-end pipeline. Raw VGRF recordings (16 sensors plus two
  per-foot totals, 100\,Hz) are trimmed, windowed and converted into token
  sequences; a bidirectional Mamba encoder produces a window score, and the
  scores of a subject's windows are averaged into one decision. Normalisation
  statistics and every hyper-parameter are estimated on training and
  validation data only, and each of the five subject-wise folds contributes a
  disjoint set of test subjects.}
  \label{fig:pipeline}
\end{figure*}

\subsection{Data}
\label{ssec:data}
We use the PhysioNet \emph{Gait in Parkinson's Disease} database, pooling its
three studies (Ga, Ju, Si). Each subject walked on level ground for about two
minutes while eight force sensors per foot sampled at 100\,Hz; the dataset
provides the 16 sensor channels plus the total force under each foot. Our
subset comprises the 100 subjects with a normal-walk recording, balanced by
design into 50 PD patients and 50 healthy controls. The course-provided
five-fold partition assigns 60 training, 20 validation and 20 test subjects
per fold, class-balanced, with test sets disjoint across folds so that pooling
the five test folds covers every subject exactly once.

\subsubsection{Pre-processing}
The first five seconds of each recording are discarded to remove gait
initiation. The remainder is cut into 5\,s windows with a 1\,s step
(500 samples per window), yielding 10\,407 windows in total, between 35 and
254 per subject. Each window inherits its subject's label. Channel
normalisation is a per-channel z-score whose mean and standard deviation are
computed over the training windows of the current fold only and then applied
unchanged to validation and test windows.

Three token representations were compared, all sharing the above windowing:
\textbf{(i) raw samples}, each window a $18 \times 500$ matrix;
\textbf{(ii) amplitude-statistic tokens}, one token per second holding 11
statistics (mean, standard deviation, skewness, kurtosis, RMS, zero-crossing
rate, median, minimum, maximum, spectral energy, spectral entropy) for each of
the 18 channels, i.e.\ 198 dimensions per token; and
\textbf{(iii) stride-cycle tokens}, one token per gait cycle carrying stride,
stance and swing time, stance ratio, normalised peak force and impulse,
stride-time increment and double-support time for both feet. Foot contacts are
detected by thresholding the per-foot total force at 5\,\% of its maximum, and
strides outside 0.5--2.5\,s are rejected as implausible; the 100 subjects give
between 34 and 166 strides each.

\subsubsection{Materials/Instruments used}
Models were trained on a Kaggle NVIDIA Tesla T4 (16\,GB, compute capability
7.5) with PyTorch 2.10 (CUDA 12.8) and \texttt{mamba-ssm} 2.3.2, using mixed
precision. Development, feature extraction and evaluation ran locally on an
Apple M-series machine. Code was version-controlled on GitHub and pulled by
the Kaggle notebook at the start of every run, so each reported number comes
from a known commit.

\subsection{Proposed approach}
\label{ssec:mod}

\subsubsection{Bidirectional Mamba classifier}
A Mamba block is a selective state-space layer: its state transition and input
projections are functions of the input, so the recurrence can forget or retain
information conditioned on content, and it runs in linear time in the sequence
length \cite{gu2023mamba}. Standard Mamba is causal, which suits language
modelling but wastes information when an entire window is available at
inference. We therefore use a bidirectional block: the normalised token
sequence is processed by a forward mixer and, after reversal, by a second
mixer with independent parameters, and the two outputs are summed before the
residual connection.

The classifier projects tokens to $d_{\text{model}}$, applies $L$ such blocks,
and pools over time by concatenating the mean and the maximum of the final
representation, which a linear layer maps to two logits. Three sizes were
studied ($d_{\text{model}}, L$) $\in \{(64,2), (128,4), (256,6)\}$; the
reported model is the largest.

\subsubsection{Training strategy}
Per fold we grid-search dropout $\in \{0.1, 0.3\}$ and learning rate
$\in \{3\!\cdot\!10^{-4}, 10^{-4}\}$, training with AdamW (weight decay
$10^{-3}$), three warm-up epochs followed by cosine decay, batch size 128 and
at most 30 epochs. The configuration with the highest validation ROC-AUC is
retrained on the union of training and validation subjects for the epoch count
that was optimal on validation, and that model predicts the test subjects --
the deep-learning analogue of \texttt{GridSearchCV(refit=True)} on a
predefined split. Early stopping uses validation ROC-AUC with patience 6 and a
minimum of 8 epochs. Because some folds occasionally converged to a degenerate
solution, a fold whose best validation AUC stays below 0.70 is retrained with
a different seed (at most three attempts). The final model averages the window
scores of three seeds.

\subsubsection{Decision rule and subject-level aggregation}
The network scores windows, but the clinical unit is the subject. A subject's
score is the mean of its window scores. The decision threshold is calibrated
per fold on that fold's \emph{validation} subjects by maximising balanced
accuracy, using the model trained on training data only; test data never
influences it. Results at the fixed threshold 0.5 are also reported and differ
by at most two points.

\subsubsection{Evaluation metrics}
We report accuracy, precision, recall and F1 (weighted by class support, so
that recall equals accuracy by construction), together with sensitivity,
specificity and ROC-AUC, plus confusion matrices. All values are the mean and
standard deviation over 1000 bootstrap resamples of the pooled test
predictions -- resampling subjects for subject-level metrics and windows for
window-level metrics.
